% This template was initially provided by Dulip Withanage.
% Modifications 
% were made by Conny Junghans and Jannik Strötgen
% and Sascha Hunold

\documentclass[
     12pt,                    % font size
     a4paper,             % paper format
     BCOR10mm,     % binding correction
     DIV14,                 % stripe size for margin calculation
     listof=totoc,                    % table listing in toc
     bibliography=totoc,       % bibliography in toc
     index=totoc,              % index in toc
%     parskip            % paragraph skip instead of paragraph indent
     twoside,
     headsepline
     ]{scrreprt}

%%%%%%%%%%%%%%%%%%%%%%%%%%%%%%%%%%%%%%%%%%%%%%%%%%%%%%%%%%%%


% PACKAGES:

% Use German :
% Input encoding
\usepackage[utf8]{inputenc}
% Font encoding
\usepackage[T1]{fontenc}
% Index-generation
\usepackage{makeidx}
% Einbinden von URLs:
\usepackage{url}
% Special \LaTex symbols (e.g. \BibTeX):
\usepackage{doc}
% Include Graphic-files:
\usepackage{graphicx}

\usepackage{setspace}

\usepackage{amssymb}
\usepackage{amsmath}

\usepackage{amsthm}
\theoremstyle{definition}
\newtheorem*{definition}{Definition}
\newtheorem*{example}{Example}
\theoremstyle{remark}
\newtheorem*{remark}{Remark}


% hyperrefs in the documents
\usepackage[bookmarks=true,colorlinks,pdfpagelabels,pdfstartview = FitH,bookmarksopen = true,bookmarksnumbered = true,linkcolor = black,plainpages = false,hypertexnames = false,citecolor = black,urlcolor=black]{hyperref} 

\usepackage[backend=biber, style=authoryear, sorting=ynt]{biblatex}
\addbibresource{references.bib}
\renewcommand*{\nameyeardelim}{\addspace}

\DeclareCiteCommand{\parencite}
  {\usebibmacro{prenote}}
  {\bibopenparen\printtext[bibhyperref]{\printnames{labelname}\nameyeardelim\printfield{year}}\bibcloseparen}
  {\multicitedelim}
  {\usebibmacro{postnote}}

\DeclareCiteCommand{\textcite}
  {\usebibmacro{prenote}}
  {\printtext[bibhyperref]{\printnames{labelname}\nameyeardelim\bibopenparen\printfield{year}\bibcloseparen}}
  {\multicitedelim}
  {\usebibmacro{postnote}}
%%%%%%%%%%%%%%%%%%%%%%%%%%%%%%%%%%%%%%%%%%%%%%%%%%%%%%%%%%%%

% OTHER SETTINGS:

% Pagestyle:
\pagestyle{headings}

% Choose language
\newcommand{\setlang}[1]{\selectlanguage{#1}\nonfrenchspacing}

\begin{document}
% TITLE:
% TITLE:
\pagenumbering{roman} 
\begin{titlepage}


\vspace*{1cm}
\begin{center}
\vspace*{3cm}
\textbf{ 
\Large Ruprecht-Karls-Universität Heidelberg\\
\smallskip
\Large Institut für Informatik\\
\smallskip
}

\vspace{3cm}

\textbf{\large Bachelor Thesis}

\vspace{0.5\baselineskip}
{\huge
Parallel ASP Planning
}
\end{center}

\vfill 

{\large
\begin{tabular}[l]{ll}
Name: & Levi Maximilian Klein\\
Student ID: & 4283957\\
Supervisor: & Dr. Gnad\\
Submission date: & 17.02.2026
\end{tabular}
}

\end{titlepage}

\onehalfspacing

\thispagestyle{empty}

\vspace*{100pt}
\noindent
Ich versichere, dass ich diese Bachelor-Arbeit selbstständig verfasst und nur die angegebenen
Quellen und Hilfsmittel verwendet habe.

\vspace*{50pt}

\noindent
Heidelberg, dd.mm.yyyy
\cleardoublepage

\section*{Zusammenfassung}

Dies ist eine Zusammenfassung der Arbeit.

\section*{Abstract}

This is the abstract.

\cleardoublepage

\tableofcontents
\cleardoublepage
\pagenumbering{arabic} 

% List of figures (Abbildungsverzeichnis):
\listoffigures
% List of tables (Tabellenverzeichnis):
\listoftables
\cleardoublepage

%%%%%%%%%%%%%%%%%%%%%%%%%%%%%%%%%%%%%%%%%%%%%%%%%%%%%%%%%%%

\chapter{Einleitung}\label{intro}
This chapter contains an overview of the topic as well as the goals and contributions of your work.

\section{Motivation}

\section{Ziele der Arbeit}

\section{Aufbau der Arbeit}
Here you describe the structure of the thesis. For example:

In Kapitel~\ref{background} werden grundlegende Methoden für diese Arbeit vorgestellt.


\chapter{Basics}\label{background}
\section{Definitions}
\begin{definition}[multivalued planning task]
    A STRIPS-like multivalued planning task according to \parencite{Helmert2006} is given by a 4 tuple $\langle \mathcal{F}, s_0, s_*, \mathcal{O} \rangle$, where:
    \begin{itemize}
      \item $\mathcal{F}$ is a finite set of state variables, also called \textbf{fluents}. Each $x \in \mathcal{F}$ has a finite domain $x^d$ and in every state, $x$ takes exactly one value from its domain.
      \item $s_0$, the \textbf{initial state}, is a total function, s.t. $s_0(x) \in x^d$ for each $x \in \mathcal{F}$
      \item $s_*$, the \textbf{goal conditions}, is a partial function, s.t. $s_*(x) \in x^d$, for each $x \in dom(s_*)$
      \item $\mathcal{O}$ is a finite set of operators, also called \textbf{actions}. We define an action as $a = \langle a^c, a^e \rangle$, where $a^c$ is the precondition of $a$ and $a^e$ is the postcondition of $a$. Both $a^c$ and $a^e$ are partial variable assignments over $\mathcal{F}$
    \end{itemize}
\end{definition}
\begin{definition}[successor state]
  Let $s$ be state and $a \in \mathcal{O}$ an action.\\
  The \textbf{successor state} $o(a,s)$, obtained by applying action $a$ in state $s$, is defined if the precondtion of $a$ holds, i.e., $a^c(x) = s(x)$ for all $x \in dom(a^c)$. Then
  $$
  s'(x) = o(a,s) = \begin{cases} a^e(x) & \text{if } x \in dom(a^e) \\ s(x) & \text{otherwise.} \end{cases}
  $$ 
\end{definition}
\begin{remark}
  Let $A_n = \langle a_1, \dots, a_n \rangle$ be a sequence of actions with length n. We define the application of $A_n$ to a state $s$ recursevily by 
  $$
    o(A_n,s)) = o(a_n,o(A_{n-1},s))
  $$
  if $o(A_i, s)$ is defined for all $1 \le i \le n$ 
\end{remark}
\begin{definition}[sequential plan]
  A \textbf{sequential plan} is a sequence of actions $A_n = \langle a_1, \dots, a_n \rangle$, s.t. 
  $s' = o(A_n,s)$ is defined and $s'(x) = s_*(x)$ for all $x \in dom(s_*)$.
\end{definition}
\printbibliography
\end{document}